\Needspace{31\baselineskip}\vspace{6mm}
\pdfbookmark[1]{Brief von Gauss an Schumacher. 1836}{letter546}
\begin{center}
{\large\itshape Brief von Gauss an Schumacher.\par}
\vspace{5mm}\hrule\vspace{2mm}
{\small \textsc{Schumacher}. Astronomische Nachrichten Nr.\ 310. 11. Juni 1836.\par}
\vspace{2mm}\hrule\vspace{5mm}
\end{center}
\begin{flushright}Göttingen. 23. April 1836.\end{flushright}
Heute Morgen zeigte die Magnetnadel außerordentliche Bewegungen, noch größere als am 7. Februar 1835. Das veranlasste mich, einige Beobachtungsreihen in der Sternwarte aufzunehmen, während Dr.\ \textsc{Goldschmidt} im magnetischen Observa\SourceLeaf{547}torium aufzeichnete. Der gleichförmige Verlauf bestätigte sich hier so schön, dass ich es wagte, den Unterschied der beiden Uhrstände daraus abzuleiten.

\Needspace{9\baselineskip}
\begin{center}\small
\begin{tabular}{@{}p{0.82\linewidth}r@{}}
Aus einem schnellen Anstieg ergab sich für den Vorsprung von Campa vor Shelton & $4'\,41''1$\\
Aus einem Abstieg wenige Minuten später & $4\quad42.4$\\\cline{2-2}
\multicolumn{1}{r}{Mittel} & $4'\,41''7$
\end{tabular}
\end{center}
Ein direkter Vergleich der Uhren ergab
\begin{center}\small
\begin{tabular}{@{}lll@{}}
1) durch ein Zeichen am Fenster & . . . . . . . . . . & $4'\,41''5$\\
2) durch einen Induktionsimpuls & . . . . . . . . . . & $4\quad41.5$
\end{tabular}
\end{center}
Also eine herrliche Bestätigung dessen, was ich in den Astronomischen Nachrichten Nr.\ 276 [S.\ 539 dieses Bandes] gesagt habe.

\clearpage
